% Esquisse d'un programme (cote 119, copy 2): archive cote 119, site page 99 (PDF page 100), typescript page 21.
% typescript. Draft compared: esquisse/pages/esq_119_023.tex (copy 1, site 23).
% AI-transcribed by Claude Opus 5.5 (claude-opus-5-5, adaptive thinking, effort high), Message Batches API, batch esq119_full (msgbatch_015xqDHkTumDysRF1qngZhRJ), 2026-11-13; not yet checked by eye.
\scanpage{119/99}
\mspage{21}

Ma réflexion sporadique sur cette question, pendant quelques années,
s'est bornée à dégager quelques principes de base simples, en at-%
tachant d'abord mon attention au cas des polyèdres réguliers
\uline{épinglés}, ce qui réduit à un minimum le bagage conceptuel néces-%
saire, et élimine pratiquement les questions de rationalité tant
soit peu délicates. Pour un tel polyèdre, on trouve une base (ou
repère) canonique de l'espace affine ou projectif ambiant, de
telle façon que les opérations du groupe cartographique $\underline{C}_n$, en-%
gendré par les réflexions fondamentales $\sigma_i$ $(0 \leq i \leq n)$, s'y
écrivent par des formules universelles, en termes de $n$ paramètres
$\alpha_1,\ldots,\alpha_n$ , qui géométriquement s'interprètent comme les doubles
des cosinus des ``angles fondamentaux'' du polyèdre. Le polyèdre se
reconstitue à partir de cette action, et du drapeau affine ou
projectif associé à la base choisie, en transformant ce drapeau
par tous les éléments du groupe engendré par les réflexions fon-%
damentales. Ainsi le n-polyèdre épinglé ``universel'' est-il défini
canoniquement sur l'anneau de polynômes à $n$ indéterminées
\[ \mathbb{Z}[\underline{\alpha}_1,\ldots,\underline{\alpha}_n] , \]
ses spécialisations sur des corps de base arbitraires $k$ (via
des valeurs $\alpha_i \in k$ données aux indéterminées $\underline{\alpha}_i$) donnant des
polyèdres réguliers correspondant à des types combinatoires di-%
vers. Dans ce jeu, il n'est pas question de se borner à des po-%
lyèdres réguliers finis, ni même à des polyèdres réguliers dont
les facettes soient d'ordre fini, i.e. pour lesquels les para-%
mètres $\alpha_i$ soient des racines d'équations ``semicyclotomiques''
convenables, exprimant que les ``angles fondamentaux'' (dans le
cas où le corps de base est $\mathbb{R}$) sont commensurables à $2\pi$. Déjà
quand $n = 1$, le polygône régulier peut-être le plus intéressant
de tous (moralement celui du polygone régulier à un seul côté !)
est celui qui correspond à $\alpha = 2$, donnant lieu à une conique cir-%
conscrite parabolique, i.e. tangente à la droite à l'infini.

Le cas fini est celui où le groupe engendré par les réflexions
fondamentales, qui est aussi le groupe des automorphismes du

\margin{99}
